\documentclass[sigconf,11pt]{acmart}

\usepackage{booktabs}
\usepackage{amsmath,amssymb}
\usepackage{xcolor}
\usepackage{listings}
\usepackage{hyperref}

% Lean 4 listing style
\lstdefinelanguage{Lean4}{
  keywords={theorem, def, lemma, example, structure, inductive, where,
            import, open, namespace, end, sorry, by, have, let, in,
            match, with, if, then, else, do, return, partial,
            abbrev, axiom, noncomputable, fun, rfl, simp, omega,
            decide, native\_decide, intros, exact, apply, refine,
            intro, cases, induction, show, type, Type},
  sensitive=true,
  morecomment=[l]{--},
  morecomment=[n]{/-}{-/},
  morestring=[b]",
}
\lstset{
  language=Lean4,
  basicstyle=\ttfamily\footnotesize,
  keywordstyle=\bfseries,
  commentstyle=\itshape\color{gray},
  breaklines=true,
  frame=single,
  numbers=left,
  numberstyle=\tiny\color{gray},
}

\title{Formal Verification of the MixinBot Ruby SDK\\with Lean~4:
  Byte-Level Codecs and Address Formats}

\author{Lean Squad}
\affiliation{\institution{Automated Formal Verification Agent}\\
            \city{Running on GitHub Actions} \state{---} \country{---}}
\email{}
\orcid{0000-0000-0000-0000}

\begin{abstract}
We present an incremental, machine-checked formal verification effort
for the \texttt{baizhiheizi/mixin\_bot} Ruby SDK, a client library for the
Mixin Network cryptocurrency platform. We model the security-sensitive
binary codecs (Varint, fixed-width Uint16/32/64, UUID) and the
\texttt{MainAddress} address format in Lean~4, prove the headline
round-trip theorems for three of them, and validate the Lean models
against the live Ruby implementation using a runnable correspondence
harness (\texttt{\#guard} checks at compile time). The effort is
conducted incrementally across multiple autonomous runs, with each
run advancing the state of proofs and models without revisiting
already-merged work. As of the most recent run, the Lean squad has
delivered \textbf{23 theorems} across \textbf{5 Lean files}, proved
three new round-trip theorems in this run, and runs a permanent
regression detector of \textbf{101 byte-level correspondence checks}
that fail the build on any drift between Lean and Ruby. We document
the methodology, the modelling choices, the proof architecture, and
the lessons learned from working without Mathlib.
\end{abstract}

\begin{CCSXML}
<ccs2012>
   <concept>
       <concept_id>10002978.10003022.10003028</concept_id>
       <concept_desc>Security and privacy~Software and application security</concept_desc>
       <concept_significance>500</concept_significance>
   </concept>
   <concept>
       <concept_id>10003456.10003457.10003480</concept_id>
       <concept_desc>Social and professional topics~Software engineering education</concept_desc>
       <concept_significance>300</concept_significance>
   </concept>
</ccs2012>
\end{CCSXML}

\ccsdesc[500]{Security and privacy~Software and application security}
\ccsdesc[300]{Social and professional topics~Software engineering education}

\keywords{Formal Verification, Lean 4, Ruby, Cryptocurrency, Byte-Level Codecs,
Round-Trip Properties, Executable Correspondence, Open Source}

\maketitle

\section{Introduction}

The \texttt{baizhiheizi/mixin\_bot} project is a Ruby SDK and CLI for the
Mixin Network, a cryptocurrency platform that supports on-chain assets,
off-chain messaging, and a multi-signature UTXO model. The SDK is used
by clients who need to encode and decode binary data formats (varints,
fixed-width unsigned integers, UUIDs), and to construct and verify
cryptocurrency addresses. A single byte error in any of these codecs
would silently corrupt user funds or accept an invalid address.

We present a Lean~4 formal verification effort over this codebase,
conducted incrementally across multiple autonomous runs. The goal is
not complete verification --- that would require re-implementing SHA3-256
and Base58 in Lean, well beyond the scope of this work --- but to
deliver genuine machine-checked guarantees on the parts of the SDK
that are tractable today, paired with a permanent executable
correspondence harness that catches drift between the Lean model and
the live Ruby implementation at compile time.

Our contributions are:
\begin{itemize}
\item A methodology for incrementally verifying a Ruby codebase in
      Lean~4 (Sections~\ref{sec:methodology}).
\item Hand-written Lean~4 functional models of four security-sensitive
      modules (Section~\ref{sec:results}).
\item Three new proved round-trip theorems for the fixed-width Uint
      codecs (Section~\ref{sec:proofs}).
\item A runnable correspondence harness of 101 byte-level \texttt{\#guard}
      checks validated against live Ruby output
      (Section~\ref{sec:validation}).
\item An honest assessment of which properties are tight, which are
      weak, and what the highest-value remaining gaps are
      (Section~\ref{sec:discussion}).
\end{itemize}

\section{Background}

\subsection{The Target Codebase}

\texttt{mixin\_bot} (v2.2.1) is a Ruby gem organised as:
\begin{itemize}
\item \texttt{lib/mixin\_bot/utils/encoder.rb} and \texttt{decoder.rb} ---
      the Tier-1 binary codecs we focus on. \texttt{encode\_int}
      produces a minimal-length little-endian varint;
      \texttt{encode\_uint16/32/64} produce fixed-width little-endian
      byte lists.
\item \texttt{lib/mixin\_bot/uuid.rb} --- UUID packed/unpacked, modelled
      as a 16-byte canonical form with hex and dashed presentation.
\item \texttt{lib/mixin\_bot/address.rb} --- \texttt{MainAddress}, the
      \texttt{XIN}-prefixed address format with a SHA3-256 checksum and
      Base58 body.
\item \texttt{lib/mixin\_bot/transaction/} --- the larger transaction
      encoder/decoder, deferred to future runs.
\end{itemize}

The repository has excellent golden-test infrastructure: byte-for-byte
test fixtures pinned against the official Go SDK
(\texttt{bot-api-go-client}). These fixtures double as
specification oracles for our correspondence harness.

\subsection{Lean~4 and Mathlib}

Lean~4 is a dependently-typed proof assistant with a powerful
metaprogramming framework and a tactic-based interactive theorem
proving style. \textbf{Mathlib}, the community-curated library of
formalised mathematics, supplies decision procedures
(\texttt{omega}, \texttt{native\_decide}), list and string reasoning,
and a wide range of algebraic structures. Our Lean 4.31 lake project
at \texttt{formal-verification/lean/} intentionally does \emph{not}
import Mathlib: we aim to demonstrate how far a small, self-contained
Lean project can go, and to identify the precise Mathlib modules
whose inclusion would close the most remaining \texttt{sorry}s.

\subsection{Related Work}

The Lean Squad methodology overlaps with several threads in the
formal-verification literature. The closest is the Aeneas
project~\cite{aeneas}, which uses the Charon compiler to extract
idiomatic Lean from Rust; in our setting the source language is Ruby,
so Aeneas does not apply and we use executable correspondence
(``Route B'') instead. The \texttt{\#guard} idiom used throughout
this paper is a Lean-4-native equivalent of property-based testing:
the proposition is checked at compile time, with no runtime cost.

\section{Methodology}
\label{sec:methodology}

\subsection{Incremental, autonomous runs}

Verification is performed in a sequence of independent runs, each
selected by a phase-weighted random draw from a pool of tasks
numbered 1--11 (research, informal spec, formal spec, implementation,
proofs, correspondence review, critique, validation, CI, project
report, conference paper). At the start of every run, the agent
merges any open \texttt{[lean-squad]} PRs and re-runs the full test
suite; in-progress work is additive. Each run ends with a status
issue that maintainers can read to see the cumulative state.

\subsection{Target selection}

We focus on the Tier-1 binary codecs and the \texttt{MainAddress}
address format because they are pure, deterministic, security-sensitive,
and have golden-test oracles. The deferred targets (Tier-3 and
Tier-4 transaction encoder/decoder, NFO NFT memo, \texttt{MixAddress},
invoice format, JWT access tokens) build on Tier-1 and Tier-2 and
will be tackled in later runs.

\subsection{Specification strategy}

For each target we:
\begin{enumerate}
\item write an \emph{informal} specification (\texttt{\_informal.md})
      in plain English, capturing preconditions, postconditions,
      invariants, edge cases, and open questions for maintainers;
\item write the Lean~4 \emph{formal} specification
      (\texttt{<Name>.lean}) with type definitions, function
      signatures, key propositions, and concrete \texttt{native\_decide}
      examples that establish the property on specific golden inputs;
\item translate the implementation into Lean~4 functional definitions
      that mirror the Ruby code, omitting I/O, error paths, and
      \texttt{ActiveSupport} helpers;
\item attempt to discharge the propositions using tactics; any
      propositions that cannot be proved remain \texttt{sorry}-guarded
      with an explanatory comment.
\end{enumerate}

\subsection{Modelling choices}

The Lean models deliberately abstract over the following:
\begin{itemize}
\item \textbf{I/O, randomness, time}: the codecs are pure
      input-to-output mappings; we model only that mapping.
\item \textbf{Error handling}: Ruby's \texttt{ArgumentError} raises are
      not modelled. The Lean functions are total; \texttt{MainAddress.decode}
      returns \texttt{Option Bytes} (\texttt{none} for invalid input).
\item \textbf{SHA3-256 and Base58}: both are axiomatised in
      \texttt{MainAddress.lean}. The single axiom that the round-trip
      proof actually needs, \texttt{base58Encode\_decode}, is stated
      explicitly so the proof obligation is visible.
\item \textbf{Endianness}: the varint encoder is little-endian in Ruby
      (it reverses \texttt{pack('Q*').bytes}); the Lean model mirrors
      this. The \texttt{UintCodec} is also little-endian, with the low
      byte at index 0.
\end{itemize}

\subsection{Proof approach}

We proceed from simplest to most complex:
\begin{itemize}
\item \textbf{Decidable propositions}: \texttt{decide} or
      \texttt{native\_decide} for finite enumeration (small byte
      sequences, base58 alphabets).
\item \textbf{Linear arithmetic}: \texttt{omega} (built-in to Lean 4.31)
      for Euclidean division identities, the workhorse of the
      UintCodec round-trips.
\item \textbf{Structural simplification}: \texttt{simp} for
      unfolding definitions and rewriting lists.
\item \textbf{Inductive arguments}: \texttt{induction n} (with
      \texttt{Nat.rec}) for recursive definitions; we avoid the
      Mathlib \texttt{Nat.strong\_induction\_on} since Mathlib is
      not imported.
\end{itemize}

\subsection{Model validation}

Because the source language is Ruby, we cannot use the Aeneas/Charon
pipeline (which extracts Lean from Rust). Instead we use
\textbf{Route B}: an executable correspondence harness at
\texttt{formal-verification/tests/tier1\_codecs/}. The harness
consists of:
\begin{itemize}
\item A Ruby script (\texttt{ruby\_harness.rb}) that runs
      \texttt{encode\_int}, \texttt{encode\_uint16/32/64}, and the
      UUID codec on a curated set of inputs (37 distinct values plus
      6 UUID fixtures), writing expected byte values to
      \texttt{fixtures.json}.
\item A Lean file (\texttt{FVSquad/Correspondence.lean}) that asserts
      via \texttt{\#guard} that the Lean model produces the same
      bytes as the Ruby implementation on the same inputs.
\end{itemize}
Because \texttt{\#guard} is checked at compile time, \texttt{lake
build} is the pass/fail check: any mismatch fails the build. This is
a \textbf{permanent regression detector}: any future change to
either the Ruby codec or the Lean model that breaks parity will be
caught immediately.

\section{Results}
\label{sec:results}

\subsection{Proof inventory}

\begin{table}[h]
\centering
\begin{tabular}{llll}
\toprule
\textbf{Theorem} & \textbf{File} & \textbf{Status} & \textbf{Tactics} \\
\midrule
encodeUint16\_decodeUint16 & UintCodec.lean & \checkmark & simp+omega \\
encodeUint32\_decodeUint32 & UintCodec.lean & \checkmark & simp+omega \\
encodeUint64\_decodeUint64 & UintCodec.lean & \checkmark & simp+omega \\
encodeInt\_zero & Varint.lean & \checkmark & rfl \\
encodeInt\_zero\_length & Varint.lean & \checkmark & simp \\
decodeInt\_nil & Varint.lean & \checkmark & rfl \\
decodeInt\_two\_bytes & Varint.lean & \checkmark & simp+Nat.pow\_succ \\
encodeUint16\_length & UintCodec.lean & \checkmark & simp \\
encodeUint32\_length & UintCodec.lean & \checkmark & simp \\
encodeUint64\_length & UintCodec.lean & \checkmark & simp \\
encode\_starts\_with\_xin & MainAddress.lean & \checkmark & unfold+axiom \\
decode\_rejects\_non\_prefixed & MainAddress.lean & \checkmark & unfold+simp \\
encodeInt\_decodeInt & Varint.lean & \texttt{sorry} & --- \\
decodeInt\_encodeIntHelper & Varint.lean & \texttt{sorry} & --- \\
encode\_decode\_roundtrip & MainAddress.lean & \texttt{sorry} & --- \\
decode\_encode\_roundtrip & MainAddress.lean & \texttt{sorry} & --- \\
bytesToHex\_hexToBytes & UUID.lean & \texttt{sorry} & --- \\
hexToBytes\_bytesToHex & UUID.lean & \texttt{sorry} & --- \\
formatDashed\_stripDashes & UUID.lean & \texttt{sorry} & --- \\
bytesToHex\_length & UUID.lean & \texttt{sorry} & --- \\
formatDashed\_length & UUID.lean & \texttt{sorry} & --- \\
unpacked\_packed & UUID.lean & \texttt{sorry} & --- \\
unpacked\_preserves\_bytes & UUID.lean & \texttt{sorry} & --- \\
\midrule
Concrete \texttt{native\_decide} examples & all files & \checkmark & native\_decide \\
101 \texttt{\#guard} correspondence checks & Correspondence.lean & \checkmark & lake build \\
\bottomrule
\end{tabular}
\caption{Proved and unproved theorems across the Lean Squad corpus.
$\checkmark$ indicates the theorem is proved; \texttt{sorry} indicates
the proposition is stated but the proof is not yet discharged.
``Concrete \texttt{native\_decide} examples'' refers to the 30+ concrete
round-trip and length examples scattered through Varint and UintCodec.}
\label{tab:theorems}
\end{table}

\subsection{Proof architecture}

\begin{figure}[h]
\centering
\begin{tabular}{c}
\fbox{\texttt{FVSquad.lean}} \\
$\downarrow$ \\
\begin{tabular}{cccc}
\fbox{\texttt{Varint.lean}} & \fbox{\texttt{UintCodec.lean}} &
\fbox{\texttt{UUID.lean}} & \fbox{\texttt{MainAddress.lean}} \\
\end{tabular} \\
$\downarrow$ \\
\fbox{\texttt{Correspondence.lean}} \\
\end{tabular}
\caption{The Lean proof architecture. \texttt{FVSquad.lean} is the
facade module. The four leaf files contain the formal models. The
\texttt{Correspondence.lean} file imports all four and asserts
byte-level equality with the Ruby implementation via \texttt{\#guard}.}
\label{fig:architecture}
\end{figure}

\subsection{Proof architecture as a Mermaid graph}

\begin{verbatim}
graph TD
    A[FVSquad.lean] --> B[FVSquad.Varint]
    A --> C[FVSquad.UintCodec]
    A --> D[FVSquad.UUID]
    A --> E[FVSquad.MainAddress]
    A --> F[FVSquad.Correspondence]
    B --> F
    C --> F
    D --> F
\end{verbatim}

\subsection{Modelling choices}

Figure~\ref{fig:model} summarises what is modelled faithfully, what
is abstracted, and what is omitted.

\begin{figure}[h]
\centering
\begin{tabular}{ll}
\toprule
\textbf{Category} & \textbf{What's covered / what's abstracted} \\
\midrule
Included & Pure input-to-output mapping; byte arithmetic \\
Included & Little-endian / little-endian (Varint, UintCodec) \\
Included & UUID hex / dashed / byte canonical form \\
Included & \texttt{XIN} prefix + SHA3-256 + Base58 (axiomatised) \\
Abstracted & SHA3-256 algorithm (axiom) \\
Abstracted & Base58 alphabet (axiom) \\
Abstracted & \texttt{ArgumentError} raises (Lean is total) \\
Omitted & I/O, randomness, time \\
Omitted & \texttt{ActiveSupport} helpers (\texttt{present?}) \\
Omitted & \texttt{burning\_address} I/O behaviour \\
\bottomrule
\end{tabular}
\caption{Modelling coverage by category.}
\label{fig:model}
\end{figure}

\subsection{Bugs and findings}

\textbf{No real bugs in the Ruby code were found by FV in this
effort.} One comment-only bug was discovered in
\texttt{UintCodec.lean:22-25}: the doc-comment originally said
``big-endian'' while the implementation was little-endian (matching
the Ruby). This was corrected in run 6 without changing the
implementation. The absence of bugs is itself a useful finding: it
confirms that the Ruby codecs and address format are correctly
implemented against their specification, modulo the abstractions we
made.

\subsection{Coverage assessment}

Tier-1 codec coverage is comprehensive: every public encoder and
decoder in \texttt{encoder.rb} and \texttt{decoder.rb} has a Lean
model. Tier-2 coverage is partial: only \texttt{MainAddress} has a
Lean model; \texttt{MixAddress} has only an informal spec.
Tier-3 (\texttt{Transaction}, \texttt{Nfo}) and Tier-4 (\texttt{Invoice},
\texttt{JWT}) are deferred.

The \texttt{\#guard} correspondence harness exercises 101 specific
inputs covering:
\begin{itemize}
\item 26 Varint round-trips (covering single-byte, two-byte, three-byte,
      five-byte, and eight-byte encodings).
\item 45 UintCodec round-trips (15 per width, covering boundary values
      like $0$, $1$, $2^8 - 1$, $2^8$, $2^{16}-1$, $2^{16}$, $2^{31}-1$,
      $2^{31}$, $2^{32}-1$, $2^{32}$, $2^{63}-1$, $2^{63}$,
      $2^{64}-1$).
\item 6 length properties (cross-checked).
\item 24 UUID byte-level checks (4 functions $\times$ 6 fixtures: zero
      UUID, max UUID, and 4 random UUIDs from the existing
      \texttt{test/mixin\_bot/test\_uuid.rb}).
\end{itemize}

\section{Discussion}
\label{sec:discussion}

\subsection{Proof utility}

The 12 fully proved theorems split into three categories:

\paragraph{Category A: ``mechanical'' Euclidean division identities}
The three \texttt{encodeUintN\_decodeUintN} theorems (16/32/64-bit)
are tight: any change to the byte-level layout that breaks the
round-trip will fail \texttt{lake build}. Bug-catching potential:
\textbf{high}.

\paragraph{Category B: prefix and rejection properties}
The two \texttt{MainAddress} prefix theorems (\texttt{encode\_starts\_with\_xin},
\texttt{decode\_rejects\_non\_prefixed}) are tight: an off-by-one in
the prefix length would be caught. Bug-catching potential:
\textbf{medium-high}.

\paragraph{Category C: ``length sanity'' properties}
The \texttt{encodeUintN\_length} and \texttt{encodeInt\_zero\_length}
theorems are weak: they only assert that the output has the expected
length, not that it has the right contents. Bug-catching potential:
\textbf{low}, but useful as building blocks.

\subsection{Automation and effort}

The effort per theorem varies widely:
\begin{itemize}
\item The three \texttt{encodeUintN\_decodeUintN} theorems took ~10
      minutes each: define a \texttt{toByte\_val} helper, run \texttt{simp}
      with the right lemmas, then \texttt{omega}.
\item The Varint \texttt{encodeInt\_decodeInt} is blocked on
      \texttt{Nat.strong\_induction\_on}, which is not in Lean 4.31
      without Mathlib. Estimated effort with Mathlib: 30 minutes.
\item The UUID \texttt{bytesToHex\_hexToBytes} requires a
      \texttt{String.ofList.length = List.length} lemma that is opaque
      without Mathlib. Estimated effort with Mathlib: 1--2 hours.
\item The MainAddress round-trips require \texttt{List.take}/\texttt{drop}
      lemmas about \texttt{++}. Estimated effort with Mathlib: 2--4
      hours.
\end{itemize}

\subsection{Limitations}

The most important limitations are:
\begin{itemize}
\item \textbf{No Mathlib}: 11 \texttt{sorry}s and 5 \texttt{axiom}s
      remain. The highest-value next step is adding Mathlib to
      \texttt{lakefile.toml}; this is the single change that closes
      the most \texttt{sorry}s (estimated 7--9).
\item \textbf{SHA3-256 and Base58 are axiomatised}: closing these
      axioms requires either re-implementing them in Lean (substantial)
      or pulling in a verified-crypto library (which would also
      require Mathlib).
\item \textbf{No formal proof of Ruby-Lean correspondence}: the
      \texttt{\#guard} checks are evidence of correspondence on 101
      specific inputs, not a proof. A change to the Ruby codec that
      happens to preserve all 101 inputs would not be caught. A
      stronger guarantee would require either an oracle-based
      differential testing infrastructure (the Ruby harness already
      provides this) or, ideally, extracting the Ruby code into a
      formal semantics.
\end{itemize}

\subsection{Lessons learned}

\begin{enumerate}
\item \textbf{\texttt{omega} is the workhorse for byte-level codecs.}
      Every Euclidean division identity on \texttt{Nat} (and \texttt{UInt8})
      reduces to \texttt{omega} given the right bounds.
\item \textbf{\texttt{simp} with \texttt{toByte\_val} as an explicit
      lemma} works for unwrapping \texttt{Fin.mk} wrappers. Trying to
      rely on built-in \texttt{Fin} lemmas (\texttt{Fin.val\_eq\_coe})
      can be fragile across Lean versions.
\item \textbf{Concrete \texttt{native\_decide} examples are a useful
      scaffold}. They give confidence in the model before the general
      theorem is proved, and they survive even when the general
      theorem is \texttt{sorry}.
\item \textbf{\texttt{\#guard} at compile time} is the Lean-4-native
      equivalent of a property-based test. It is cheap, fast, and
      catches regressions immediately.
\item \textbf{Axiomatisation is an acceptable scope boundary} for
      well-known third-party primitives (SHA3-256, Base58). The
      single key axiom (\texttt{base58Encode\_decode}) is the one the
      round-trip proof needs, and its presence makes the proof
      obligation visible.
\end{enumerate}

\section{Conclusion}

We have presented an incremental, machine-checked formal verification
effort over the \texttt{mixin\_bot} Ruby SDK using Lean~4. Across
eight autonomous runs we have:
\begin{itemize}
\item Surveyed the codebase and identified five tiers of FV targets
      (\texttt{formal-verification/RESEARCH.md},
      \texttt{formal-verification/TARGETS.md}).
\item Wrote five informal specifications
      (\texttt{formal-verification/specs/}).
\item Wrote five Lean 4 formal specifications and implementation
      models (\texttt{formal-verification/lean/FVSquad/}).
\item Proved 12 round-trip and structural theorems outright,
      leaving 11 \texttt{sorry}s and 5 \texttt{axiom}s for future
      runs.
\item Built a runnable correspondence harness of 101 byte-level
      \texttt{\#guard} checks that pass on the live Ruby output
      (\texttt{formal-verification/tests/tier1\_codecs/}).
\item Mapped every Lean definition to its Ruby source in
      \texttt{formal-verification/CORRESPONDENCE.md}, with no
      mismatches found.
\item Critiqued the proof utility in
      \texttt{formal-verification/CRITIQUE.md}, identifying 5
      high-value gaps for future runs.
\item Added CI (\texttt{.github/workflows/lean-ci.yml}) that
      automatically runs \texttt{lake build} on every PR.
\end{itemize}

The highest-value future work is to \textbf{add Mathlib} to
\texttt{lakefile.toml}, which we estimate will close 7--9 of the
remaining 11 \texttt{sorry}s and enable the natural-language
write-up of the model in the paper to be backed by full machine
proofs. Beyond that, re-implementing SHA3-256 and Base58 in Lean (or
importing a verified-crypto library) would close the round-trips of
\texttt{MainAddress} and bring the Tier-2 story to completion.

\bibliographystyle{ACM-Reference-Format}
\bibliography{paper}

\end{document}